\documentclass[dvipsnames]{book}

\usepackage[margin=1cm,a5paper]{geometry}

\usepackage{bold-extra}
\usepackage{polski}
\usepackage{iwona}
\usepackage[OT4]{fontenc}

\usepackage{amsmath,amssymb,amsthm}
\usepackage{pgf,tikz}
\usepackage{hyperref}

\pagestyle{empty}
\setlength{\parindent}{0pt}
\setlength{\parskip}{5pt}
\renewcommand\leq\leqslant
\renewcommand\geq\geqslant

\colorlet{maincolor}{ForestGreen}
\usepackage[most]{tcolorbox}
\usepackage{varwidth}

\newtcolorbox{tresc}[3][]{enhanced, sharp corners, boxrule=0mm, colback=maincolor!10,
colframe=black,toptitle=1mm,fonttitle=\bfseries,
colbacktitle=maincolor!25, coltitle=black,
title={Zadanie #2\hfill{\it Opracowanie: #3}},#1}

\newenvironment{problem}[3]{\newpage\begin{tresc}{#1}{#2}{#3}}{\end{tresc}}

% Dodatkowe makra

\usepackage{multicol}
\usepackage{cancel}
\usepackage{bbm}

\def\ind{\mathbbm{1}}

\def\N{\mathbb{N}}
\def\P{\mathbb{P}}
\def\E{\mathbb{E}}
\def\Z{\mathbb{Z}}
\def\Q{\mathbb{Q}}
\def\R{\mathbb{R}}

\def\nwd{\mathrm{NWD}}

\newtheorem*{lemat}{Lemat}

\usepackage{enumitem}
\setlist[enumerate,1]{label={(\alph*)},topsep=0pt}

% ------------------------------------------------------------------
% Kolorowe boxy: definicje (zielone), twierdzenia/obserwacje/wnioski (niebieskie),
% algorytmy (liliowe)
% ------------------------------------------------------------------

\colorlet{defcolor}{ForestGreen}
\colorlet{thmcolor}{RoyalBlue}
\colorlet{algcolor}{Orchid}

\newcounter{defcnt}
\newcounter{thmcnt}
\newcounter{algcnt}

% wspólny styl boxa
\tcbset{notatkibox/.style={enhanced, sharp corners, boxrule=0mm, frame hidden,
  toptitle=0.5mm, bottomtitle=0.5mm, left=2mm, right=2mm, top=1.5mm, bottom=1.5mm,
  fonttitle=\bfseries, coltitle=black, breakable, before skip=6pt, after skip=6pt}}

% \begin{definicja}[nazwa]
\NewDocumentEnvironment{definicja}{o}{%
  \refstepcounter{defcnt}%
  \begin{tcolorbox}[notatkibox, colback=defcolor!8, colframe=defcolor!25,
    colbacktitle=defcolor!25,
    title={Definicja \thedefcnt\IfValueT{#1}{\ (#1)}}]%
}{\end{tcolorbox}}

% generator środowisk niebieskich: #1 = nazwa środowiska, #2 = etykieta
\newcommand{\newthmbox}[2]{%
  \NewDocumentEnvironment{#1}{o}{%
    \refstepcounter{thmcnt}%
    \begin{tcolorbox}[notatkibox, colback=thmcolor!8, colframe=thmcolor!25,
      colbacktitle=thmcolor!25,
      title={#2 \thethmcnt\IfValueT{##1}{\ (##1)}}]%
  }{\end{tcolorbox}}%
}

\newthmbox{twierdzenie}{Twierdzenie}
\newthmbox{obserwacja}{Obserwacja}
\newthmbox{wniosek}{Wniosek}
\newthmbox{stwierdzenie}{Stwierdzenie}

% \begin{algorytm}[nazwa] -- liliowy box
\NewDocumentEnvironment{algorytm}{o}{%
  \refstepcounter{algcnt}%
  \begin{tcolorbox}[notatkibox, colback=algcolor!8, colframe=algcolor!25,
    colbacktitle=algcolor!25,
    title={Algorytm \thealgcnt\IfValueT{#1}{\ (#1)}}]%
}{\end{tcolorbox}}

% uwaga -- żółty box, \begin{uwaga}[nazwa]
\colorlet{uwagacolor}{Dandelion}

\NewDocumentEnvironment{uwaga}{o}{%
  \refstepcounter{thmcnt}%
  \begin{tcolorbox}[notatkibox, colback=uwagacolor!15, colframe=uwagacolor!40,
    colbacktitle=uwagacolor!40,
    title={Uwaga \thethmcnt\IfValueT{#1}{\ (#1)}}]%
}{\end{tcolorbox}}

% treść zadania -- szary box, \begin{zadanie}{1.1}
\colorlet{zadcolor}{Gray}

\NewDocumentEnvironment{zadanie}{m}{%
  \begin{tcolorbox}[notatkibox, colback=zadcolor!15, colframe=zadcolor!40,
    colbacktitle=zadcolor!40,
    title={Zadanie #1}]%
}{\end{tcolorbox}}

% przykład -- bez boxa, wyróżniony nagłówkiem
\theoremstyle{definition}
\newtheorem*{przyklad}{Przykład}

\renewcommand{\qedsymbol}{$\square$}


% numerowany lemat w niebieskim boxie (wspólny licznik z twierdzeniami)
\newthmbox{lematbox}{Lemat}

\begin{document}

\begin{center}
  {\Large\bfseries Optymalizacja stochastyczna --- Wykład 2 --- 2026-10-12}
\end{center}

\section*{Teoria rekordów (cd.)}

Przypomnienie: $X_1, X_2, \ldots$ --- iid o ciągłej dystrybuancie $F$,
$N(1) = 1$, $N(k) = \min\big\{j > N(k-1) : X_j > X_{N(k-1)}\big\}$ --- indeks
$k$-tego rekordu.

\textbf{Lemat 1.} Jeśli $F$ jest ciągła, to rozkład $N(k)$ nie zależy od $F$.
(Dowód będzie na ćwiczeniach.)

Dzięki Lematowi 1 we wszystkich dowodach poniżej możemy bez straty ogólności
(BSO) przyjąć, że $\bar X_i \sim \mathcal{U}(0,1)$ iid.

% numeracja zgodna z wykładem: Lemat 1 był na poprzednim wykładzie
\setcounter{thmcnt}{1}

\begin{twierdzenie}
Rozkład $N(2)$ zadany jest przez
\[
  \P\big(N(2) = n\big) = \frac{1}{n(n-1)}, \qquad n \geq 2.
\]
\end{twierdzenie}

\begin{proof}
Niech $n \geq 2$. BSO $\bar X_i \sim \mathcal{U}(0,1)$ iid. Ze wzoru na
prawdopodobieństwo całkowite:
\begin{align*}
  \P\big(N(2) = n\big)
  &= \int_0^1 \P\big(N(2) = n \mid \bar X_1 = x\big)\, dx
  \overset{(*)}{=} \int_0^1 (1-x)\, x^{n-2}\, dx \\
  &= \frac{x^{n-1}}{n-1}\bigg|_0^1 - \frac{x^n}{n}\bigg|_0^1
  = \frac{1}{n-1} - \frac{1}{n} = \frac{1}{n(n-1)}.
\end{align*}
Ad $(*)$: zdarzenie $\{N(2) = n\}$ przy $\bar X_1 = x$ oznacza, że
\[
  \bar X_2 \in (0,x),\ \bar X_3 \in (0,x),\ \ldots,\ \bar X_{n-1} \in (0,x),
  \quad \bar X_n \in (x,1),
\]
bo w $n$-tym kroku był drugi rekord, więc ta wartość przekroczyła $x$.
Z niezależności dostajemy $x^{n-2}(1-x)$.
\end{proof}

\begin{wniosek}
W szczególności $\E\big[N(k)\big] = +\infty$ dla $k \geq 2$.
\end{wniosek}
\begin{proof}
\[
  \E\big[N(2)\big] = \sum_{n \geq 2} n \cdot \frac{1}{n(n-1)}
  = \sum_{n \geq 2} \frac{1}{n-1} = +\infty,
\]
a ponieważ $N(k) \geq N(2)$ dla $k \geq 2$, to
$\E\big[N(k)\big] \geq \E\big[N(2)\big] = +\infty$.
\end{proof}

\hrulefill

\begin{twierdzenie}
Proces $\big(N(k)\big)_{k \geq 2}$ ma własność Markowa, tzn.
\[
  \P\big(N(k) = n \mid N(l) = j_l,\ l = 2, 3, \ldots, k-1\big)
  = \P\big(N(k) = n \mid N(k-1) = j_{k-1}\big)
\]
dla każdego wektora $(j_2 < j_3 < \ldots < j_{k-1})$. Ponadto
prawdopodobieństwa przejścia spełniają
\[
  \P\big(N(k) = n \mid N(k-1) = j\big) =
  \begin{cases}
    \dfrac{j}{n(n-1)} & : n > j \geq k-1 \geq 2, \\[2mm]
    0 & : \text{wpp.}
  \end{cases}
\]
\end{twierdzenie}

\begin{proof}
BSO $\bar X_i \sim \mathcal{U}(0,1)$ iid. Oznaczmy $j_1 := 1 = N(1)$. Ze wzoru na
prawdopodobieństwo całkowite (tak samo jak w poprzednim dowodzie, tylko po
wszystkich współrzędnych $x_l$):
\begin{align*}
  &\P\big(N(l) = j_l : 2 \leq l \leq k\big) = \\
  &\quad = \int_0^1 \!\!\cdots\! \int_0^1
    \P\big(N(l) = j_l : 2 \leq l \leq k \;\big|\;
    \bar X_{j_l} = x_l : 1 \leq l < k\big)\, dx_1 \ldots dx_{k-1}
  \overset{(\heartsuit)}{=} \ldots
\end{align*}

\paragraph{Komentarz.}
Prawdopodobieństwo pod całką wynosi $0$, chyba że
$x_1 < x_2 < \ldots < x_{k-1}$ (wtedy te rekordy przyjmują właśnie te
wartości). Ponadto zdarzenie oznacza, że
\begin{align*}
  \bar X_{j_1} &= x_1, & \bar X_{j_1+1}, \ldots, \bar X_{j_2-1} &\in (0, x_1), \\
  \bar X_{j_2} &= x_2, & \bar X_{j_2+1}, \ldots, \bar X_{j_3-1} &\in (0, x_2), \\
  &\ \ \vdots & &\ \ \vdots \\
  \bar X_{j_{k-1}} &= x_{k-1}, & \bar X_{j_{k-1}+1}, \ldots, \bar X_{j_k-1} &\in (0, x_{k-1}), \\
  & & \bar X_{j_k} &\in (x_{k-1}, 1).
\end{align*}
Zatem $l$-ta grupa daje czynnik $x_l^{\,j_{l+1} - j_l - 1}$ (np.\ pierwsza
$x_1^{\,j_2 - 2}$, druga $x_2^{\,j_3 - j_2 - 1}$), a tylko ostatni czynnik jest
inny: $(1 - x_{k-1})$.

\paragraph{Koniec komentarza.}
\begin{align*}
  (\heartsuit)
  &= \int\limits_{0 < x_1 < x_2 < \ldots < x_{k-1} < 1}
    x_1^{\,j_2 - j_1 - 1} x_2^{\,j_3 - j_2 - 1} \cdots
    x_{k-1}^{\,j_k - j_{k-1} - 1} (1 - x_{k-1})\, dx_1 \ldots dx_{k-1}.
\end{align*}
Całkujemy po kolei po $x_1, x_2, \ldots$ Po $x_1$:
\[
  \int_0^{x_2} x_1^{\,j_2 - j_1 - 1}\, dx_1 = \frac{x_2^{\,j_2 - j_1}}{j_2 - j_1},
  \qquad\text{więc}\qquad
  \frac{x_2^{\,j_2 - j_1}}{j_2 - j_1} \cdot x_2^{\,j_3 - j_2 - 1}
  = \frac{1}{j_2 - j_1}\, x_2^{\,j_3 - j_1 - 1}.
\]
Po scałkowaniu po $x_2$ mamy
\[
  \int_0^{x_3} x_2^{\,j_3 - j_1 - 1}\, dx_2 = \frac{x_3^{\,j_3 - j_1}}{j_3 - j_1}
\]
itd. Po scałkowaniu po $x_1, \ldots, x_{k-2}$ zostaje
\[
  (\heartsuit) = \frac{1}{(j_2 - j_1)(j_3 - j_1) \cdots (j_{k-1} - j_1)}
  \int_0^1 x_{k-1}^{\,j_k - j_1 - 1} (1 - x_{k-1})\, dx_{k-1}.
\]
Z dowodu poprzedniego twierdzenia (bo $j_k - j_1 - 1 = j_k - 2$) ostatnia
całka wynosi $\frac{1}{j_k (j_k - 1)}$, więc
\[
  (\heartsuit) = \frac{1}{j_k \prod_{l=2}^{k} (j_l - 1)}.
\]

Podsumowując (i pamiętając, że $j_1 = 1$):
\begin{align*}
  \P\big(N(l) = j_l : 2 \leq l \leq k\big)
    &= \frac{1}{\prod_{l=2}^{k} (j_l - 1)} \cdot \frac{1}{j_k}, \\
  \P\big(N(l) = j_l : 2 \leq l \leq k-1\big)
    &= \frac{1}{\prod_{l=2}^{k-1} (j_l - 1)} \cdot \frac{1}{j_{k-1}}.
\end{align*}
Stąd lewa strona tezy:
\[
  \text{LHS}
  = \frac{\frac{1}{j_k - 1} \cdot \frac{1}{j_k}}{\frac{1}{j_{k-1}}}
  = \frac{j_{k-1}}{j_k (j_k - 1)},
\]
co nie zależy od $j_2, j_3, \ldots, j_{k-2}$. Sumując po tych indeksach
dostajemy, że $\P\big(N(k) = j_k \mid N(k-1) = j_{k-1}\big)$ ma tę samą
wartość, zatem LHS $=$ RHS. Dla $j_k \leq j_{k-1}$ prawdopodobieństwo jest
oczywiście zerowe, a $j_{k-1} \geq k-1$, bo $N(l) \geq l$.
\end{proof}

\begin{wniosek}
Rozkład $N(k)$ dla dowolnego $k \geq 2$ można wyznaczyć rekurencyjnie
za pomocą $N(l)$, $l < k$.
\end{wniosek}

\begin{przyklad}
\begin{align*}
  \P\big(N(3) = n\big)
  &= \sum_{j=2}^{\infty} \P\big(N(3) = n \mid N(2) = j\big)\,
     \P\big(N(2) = j\big) \\
  &= \sum_{j=2}^{n-1} \frac{j}{n(n-1)} \cdot \frac{1}{j(j-1)}
   = \frac{1}{n(n-1)} \sum_{j=2}^{n-1} \frac{1}{j-1}.
\end{align*}
\end{przyklad}

\hrulefill

\section*{Wskaźniki rekordów}

Rozważmy zmienne losowe $(e_n)_{n \geq 1}$ zadane następująco:
\[
  e_n = \begin{cases}
    1 & : \exists k\ N(k) = n, \\
    0 & : \text{wpp.}
  \end{cases}
\]
(czyli $e_n = 1$ $\iff$ w $n$-tym kroku padł rekord; w szczególności $e_1 = 1$).

\begin{lematbox}
Zmienne $(e_n)_{n \geq 2}$ są niezależne \textup{(1)} oraz
$\P(e_n = 1) = \frac{1}{n}$ \textup{(2)}.
\end{lematbox}

\begin{proof}
BSO $(\bar X_j)_{j \geq 1} \sim \mathcal{U}(0,1)$ iid.

\textbf{Ad (2).} Rozważmy zdarzenie $\{\omega \in \Omega : e_n(\omega) = 1\}$,
tzn.\ $\bar X_n > \bar X_j$ dla $j = 1, \ldots, n-1$. Wtedy
\[
  \P(e_n = 1)
  = \int_0^1 \P\big(\bar X_j < x : j = 1, \ldots, n-1 \mid \bar X_n = x\big)\, dx
  = \int_0^1 x^{n-1}\, dx = \frac{1}{n}.
\]

\textbf{Ad (1).} Rozważmy dowolne całkowite $2 \leq k_1 < k_2 < \ldots < k_m$.
Wtedy zdarzenie $\{e_{k_l} = 1 : l = 1, \ldots, m\}$ oznacza, że
$\forall l$: $\bar X_j < \bar X_{k_l}$ dla $j < k_l$. Zatem
\begin{align*}
  &\P\big(e_{k_l} = 1 : 1 \leq l \leq m\big) = \\
  &= \int_0^1 \!\!\cdots\! \int_0^1
    \P\big(e_{k_l} = 1 : 1 \leq l \leq m \;\big|\;
    \bar X_{k_l} = x_l : 1 \leq l \leq m\big)\, dx_1 \ldots dx_m \\
  &= \int\limits_{0 < x_1 < x_2 < \ldots < x_m < 1}
    x_1^{\,k_1 - 1} x_2^{\,k_2 - k_1 - 1} \cdots x_m^{\,k_m - k_{m-1} - 1}\,
    dx_1 \ldots dx_m \\
  &= \frac{1}{k_1 k_2 \cdots k_m}
   = \prod_{l=1}^{m} \P\big(e_{k_l} = 1\big),
\end{align*}
(całkujemy kolejno po $x_1, x_2, \ldots$ jak w dowodzie Twierdzenia 3), czyli
zachodzi reguła mnożenia prawdopodobieństw.

$\sigma$-ciało generowane przez $e_n$ to
$\{\emptyset, \Omega, \{e_n = 1\}, \{e_n = 0\}\}$, zatem jeśli reguła mnożenia
zachodzi dla zdarzeń $\{e_n = 1\}$, to dla innych zbiorów też (przechodząc do
dopełnień). Stąd niezależność.
\end{proof}

\begin{wniosek}
Przy powyższych oznaczeniach:
\begin{enumerate}
  \item $\displaystyle
    \P\Big(\lim_{n \to \infty} \frac{\sum_{k=1}^{n} e_k}{\ln n} = 1\Big) = 1$
    \hfill (coś jak prawo wielkich liczb),
  \item $\displaystyle
    \P\bigg(\limsup_{n \to \infty}
      \frac{\sum_{k=1}^{n} e_k - \ln n}{\big(2 \ln n \cdot \ln\ln\ln n\big)^{1/2}}
      = 1\bigg) = 1$
    \hfill (coś jak prawo iterowanego logarytmu),
  \item $\displaystyle
    \lim_{n \to \infty} \P\Big(\sum_{k=1}^{n} e_k - \ln n < x \sqrt{\ln n}\Big)
    = \Phi(x)$
    \hfill (coś jak CTG).
\end{enumerate}
\end{wniosek}

\begin{proof}[Dowód --- szkic]
\[
  \E[e_k] = \frac{1}{k}, \qquad
  \mathrm{Var}(e_k) = \frac{1}{k} - \frac{1}{k^2} = \frac{1}{k}\Big(1 - \frac{1}{k}\Big), \qquad
  \sum_{k=1}^{n} \frac{1}{k} \approx \ln n,
\]
więc również $\sum_{k=1}^n \mathrm{Var}(e_k) \approx \ln n$.
\begin{enumerate}
  \item mocne prawo wielkich liczb,
  \item prawo iterowanego logarytmu,
  \item CTG.
\end{enumerate}
Jedyne potrzebne założenia to $e_n \in L^\infty$ i niezależność.
\end{proof}

\hrulefill

\begin{twierdzenie}
\[
  \lim_{k \to \infty} \frac{\ln N(k)}{k} = 1
  \quad \text{z prawdopodobieństwem } 1.
\]
\end{twierdzenie}

\begin{proof}
Praca domowa.

\textbf{Wskazówka:} Lemat 5 $+$ Wniosek 6 $\implies$
$\forall \varepsilon > 0\ \exists k_0\ \forall k \geq k_0$:
\[
  (1 - \varepsilon) k \leq \ln N(k) \leq (1 + \varepsilon) k
  \quad \text{z prawdopodobieństwem} \geq 1 - \varepsilon.
\]
W tym celu należy użyć
\[
  \P\big(N(k) > m\big) = \P\big(e_1 + e_2 + \ldots + e_m < k\big)
\]
($k$-ty rekord pada po chwili $m$ $\iff$ do chwili $m$ padło mniej niż $k$
rekordów).
\end{proof}

\end{document}
